% 387_Abweichungen_En_ch.tex
% Johann Pascher, 2 Oct 2026
% The deviations at a glance: ratios, SI translation, correction quantities

\providecommand{\BM}{\textbf{[B]}}
\providecommand{\KM}{\textbf{[K]}}
\providecommand{\SM}{\textbf{[S]}}
\providecommand{\XM}{\textbf{[X]}}
\providecommand{\QM}{\textbf{[Q]}}

\chapter*{The deviations at a glance}
\addcontentsline{toc}{chapter}{The deviations at a glance}

\begin{abstract}
This document records all residual deviations of FFGFT in one place, each with the measured value
and its uncertainty. It follows the order of the theory: first the ratios of the basic form in
natural units with $\alpha=1$, then the translation into SI. The ratios agree to within per mille
to a few per cent; the measurement uncertainties are almost everywhere orders of magnitude smaller
than the deviations and contribute only a minimal amount. Three correction quantities carry the
residuals: the Higgs correction $\beta_T$, which already appeared in the earliest documents of 2025
as a small deviation from $\beta_T=1$, the fractal correction $K_{\text{frak}}=74/75$ of the SI
translation, and the factor~10 in $v/E_P$. For each one, the document shows which variants exist,
how well they fit and which connections arise.
\end{abstract}

% ============================================================
\section*{Status markers}
% ============================================================
\begin{center}
\begin{tabular}{@{}lp{11cm}@{}}
\textbf{[B]} & Algebraically proven. \\[2pt]
\textbf{[K]} & Numerically confirmed (check script, comparison with measured values). \\[2pt]
\textbf{[S]} & Assumption, open or speculative. \\[2pt]
\textbf{[X]} & Refuted or not sustainable. \\[2pt]
\textbf{[Q]} & Convention. \\
\end{tabular}
\end{center}
Measured values: CODATA 2022, PDG 2024, neutrino data as in Doc.~340 with the uncertainties from
PDG 2024. The column ``measurement error'' gives the relative uncertainty of the comparison value;
where it is comparable to the deviation, the deviation is also given in standard deviations.
Check script: \texttt{pruef\_387\_abweichungen.py}.

% ============================================================
\section{How to read the deviations}
% ============================================================

FFGFT is first and foremost ratio-based (Docs.~384, 386). In the basic form, in natural units with
$\alpha=1$ and a redefined charge, there is only $\xi=1/7500$ and chosen integers; ratios test the
structure. Absolute values in SI need an anchor and a reference energy and test the conversion
(A130). The deviations are therefore listed separately. A deviation much larger than the
measurement error is a genuine residual of the theory, not a measurement problem; this holds for
almost all entries.

% ============================================================
\section{Ratios of the basic form}
% ============================================================

\begin{center}
\small
\begin{tabular}{@{}p{3.4cm}p{3.0cm}p{2.6cm}p{2.0cm}p{2.4cm}@{}}
\toprule
\textbf{Quantity} & \textbf{FFGFT} & \textbf{Deviation} & \textbf{Meas.\ error} & \textbf{Status} \\
\midrule
$m_\mu/m_e$ & $120\sqrt3=207.85$ & $+0.52\,\%$ & $2\times10^{-8}$ & \KM, genuine residual \\
$m_\tau/m_\mu$ & $\tfrac{125}{144}\xi^{-1/3}=16.99$ & $+1.03\,\%$ & $5\times10^{-5}$ & \KM, genuine residual \\
Koide $Q-\tfrac23$ & $0$ & $-0.017\,\xi$ & $0.038\,\xi$ & $0.4\sigma$ \\
$v/E_P$ & $\xi^4/(5\pi)$ & $-0.23\,\%$ & $1\times10^{-5}$ & \KM, factor~10 \\
$\xi$ from $m_h/v$ & $\beta_T=0.977$ & $-2.3\,\%$ & $1.8\times10^{-3}$ & check \KM \\
$\Delta m^2_{\text{atm}}/\Delta m^2_{\text{sol}}$ & $120/(11/3)=32.7$ & $-1.4\,\%$ & $2.6\,\%$ & $0.5\sigma$ \\
$\sin^2\theta_{13}$ & $(25\xi)^{2/3}=0.0223$ & $+1.4\,\%$ & $3\,\%$ & $0.4\sigma$ \\
$M_W/M_Z$ & $\sqrt{7/9}$ & $+0.06\,\%$ & $1.7\times10^{-4}$ & $3.8\sigma$ \\
$m_h/M_Z$ & $11/8$ & $+0.15\,\%$ & $9\times10^{-4}$ & $1.7\sigma$ \\
$m_t/M_Z$ & $(11/8)^2$ & $-0.10\,\%$ & $1.7\times10^{-3}$ & $0.6\sigma$ \\
$\lambda_{\text{CKM}}$ & $\xi^{1/6}=0.2260$ & $+0.46\,\%$ & $3\times10^{-3}$ & $1.5\sigma$ \\
\bottomrule
\end{tabular}
\end{center}

The lepton ratios are the sharpest tests of the basic form. Their measurement errors are $10^{-8}$
and $10^{-5}$, the deviations $0.5$ and $1\,\%$; the residuals are thus entirely residuals of the
ladder, and the measurement contributes nothing \KM. Their cause is open (R113). For the neutrinos,
$\theta_{13}$, $m_t$ and $\lambda_{\text{CKM}}$ the deviations lie within 1.5 standard
deviations; there the measurement is the limiting factor. $M_W$ from $\sin^2\theta_W=2/9$, at
$3.8\sigma$, is the only stated tension in the basic form. Koide agrees within the measurement
precision of $m_\tau$.

% ============================================================
\section{The translation into SI}
% ============================================================

\begin{center}
\small
\begin{tabular}{@{}p{3.6cm}p{3.0cm}p{2.4cm}p{2.0cm}p{2.4cm}@{}}
\toprule
\textbf{Quantity} & \textbf{FFGFT} & \textbf{Deviation} & \textbf{Meas.\ error} & \textbf{Status} \\
\midrule
$\alpha^{-1}$ with $74/75$ & $137.059$ & $+1.7\times10^{-4}$ & $1.5\times10^{-10}$ & \KM, anchor \\
$\alpha^{-1}$ Galois & $3700/27=137.037$ & $+7.6\times10^{-6}$ & $1.5\times10^{-10}$ & \KM \\
$m_em_\mu$ against $54$\,MeV$^2$ & $54$ & $-1.6\times10^{-4}$ & $2\times10^{-8}$ & \KM \\
$m_e$ (chain, measured $E_P$) & $\tfrac{4}{15\pi}\xi^{11/2}E_P$ & $-1.32\,\%$ & $1\times10^{-5}$ & \KM \\
$m_\mu$ (chain) & $\tfrac{16}{25\pi}\xi^{5}E_P$ & $-0.80\,\%$ & $1\times10^{-5}$ & \KM \\
$m_\tau$ (chain) & $\tfrac{5}{9\pi}\xi^{14/3}E_P$ & $+0.22\,\%$ & $5\times10^{-5}$ & \KM \\
$v$ from the ladder, bare & $248.93$\,GeV & $+1.10\,\%$ & $10^{-6}$ & \KM \\
$v$ from the ladder with $K$ & $245.61$\,GeV & $-0.25\,\%$ & $10^{-6}$ & \KM \\
$\Delta m^2_{\text{atm}}$ & $2.476\times10^{-3}$\,eV$^2$ & $-1.0\,\%$ & $1.1\,\%$ & $0.9\sigma$ \\
$\Delta m^2_{\text{sol}}$ & $7.565\times10^{-5}$\,eV$^2$ & $+0.46\,\%$ & $2.4\,\%$ & $0.2\sigma$ \\
\bottomrule
\end{tabular}
\end{center}

Here too the measurement errors, except for the neutrinos, are orders of magnitude smaller than the
deviations. Two residuals stand out. The residual of $\alpha$, $1.7\times10^{-4}$, is exactly the
distance of the measured masses from the Galois value: $m_em_\mu=53.991$\,MeV$^2$ lies
$1.6\times10^{-4}$ below $54=|\mathrm{GF}(3)^*|\cdot|\mathrm{GF}(27)|$, and with $54$,
$\alpha^{-1}=3700/27$ agrees to $7.6$\,ppm \KM. The residual of the electron mass in the chain,
$-1.32\,\%$, equals $K_{\text{frak}}-1=-1.33\,\%$ to within $1.6\times10^{-4}$ \KM; this does not
hold for the muon and the tau.

% ============================================================
\section{The correction quantities}
% ============================================================

\subsection{\texorpdfstring{$\beta_T$}{beta\_T}: the Higgs correction from the earliest documents}

The earliest documents of March and April 2025 set two quantities to~1: $\alpha=1$ and
$\beta_T=1$, on the grounds that both are conventions in natural units and that their SI values,
$\alpha\approx1/137$ and $\beta\approx0.008$, arise only through the translation (``The consistency
of $\alpha=1$ and $\beta=1$'', version history). $\beta_T=1$ with $r_0=\xi\,\ell_P$ gave the Higgs
formula $\xi=\lambda_h^2v^2/(16\pi^3m_h^2)$; the document gave $r_0\approx\ell_P/7519$ instead of
$\ell_P/7500$, i.e. already then a small deviation from $\beta_T=1$ (Doc.~385). Computed today,
\begin{equation}
  \beta_T=\frac{\xi_{\text{EFT}}}{\xi}=\frac{m_h^2}{64\pi^3v^2\,\xi}=0.977\quad\KM ,
\end{equation}
with a measurement error of $0.18\,\%$ from $m_h$. The deviation from $\beta_T=1$ is therefore a
genuine residual, 13 times larger than the measurement uncertainty. It depends on which $v$ is
inserted:
\begin{center}
\small
\begin{tabular}{@{}lcc@{}}
\toprule
\textbf{Input for $v$} & $v$ & $\beta_T$ \\
\midrule
measured (from $G_F$) & $246.22$\,GeV & $0.977$ \\
chain $\xi^4E_P/(5\pi)$ & $245.65$\,GeV & $0.982$ \\
ladder with $m_e$ and $K_{\text{frak}}$ & $245.61$\,GeV & $0.982$ \\
ladder with $m_e$, bare & $248.93$\,GeV & $0.956$ \\
\bottomrule
\end{tabular}
\end{center}
With FFGFT's own values for $v$, $\beta_T$ moves closer to~1 but stays at $-1.8\,\%$; with the bare
$v$ it moves away. The Higgs check remains a check of $\xi$ to within a few per cent \KM. There is
no derivation of the residual \SM. Numerically, $\beta_T=0.9772$ lies close to
$K_{\text{frak}}^{\sqrt3}=0.9770$; with a measurement error of $0.18\,\%$ this is not a robust
statement.

\subsection{\texorpdfstring{$K_{\text{frak}}$}{K\_frak}: the correction of the SI translation}

In the basic form there is no correction factor; the bridge reads $\xi E_0^2=1$ (Doc.~386). In the
SI translation, $\alpha$ requires the factor $K_\alpha=\xi m_em_\mu\alpha^{-1}=0.98650$. The variants
compared:
\begin{center}
\small
\begin{tabular}{@{}lccccc@{}}
\toprule
\textbf{Variant} & $K$ & $\alpha^{-1}$ & \textbf{Galois} & $v$ (ladder) & \textbf{Turns} \\
\midrule
$1-100\xi=74/75$ & $0.986667$ & $+1.7\times10^{-4}$ & $+7.6$\,ppm & $-0.25\,\%$ & $75.0$ \\
$(1-\xi)^{100}\approx e^{-100\xi}$ & $0.986754$ & $+2.6\times10^{-4}$ & $+96$\,ppm & $-0.24\,\%$ & $75.5$ \\
$1/(1+100\xi)$ & $0.986842$ & $+3.5\times10^{-4}$ & $+190$\,ppm & $-0.23\,\%$ & $76.0$ \\
$1-\psi'(75)$ (Doc.~381) & $0.986577$ & $+7.8\times10^{-5}$ & $-83$\,ppm & $-0.26\,\%$ & $74.5$ \\
without correction & $1$ & $+1.4\,\%$ & $+1.4\,\%$ & $+1.10\,\%$ & -- \\
\bottomrule
\end{tabular}
\end{center}
Only $74/75$ closes the winding after a whole number of turns and matches the Galois form to
$7.6$\,ppm \KM. $1-\psi'(75)$ halves the $\alpha$ residual but loses the whole number of turns and
worsens the Galois form. The multiplicative and the inverse forms fit worse everywhere than the
additive one; regarding the open form (R139), the $\alpha$ sector thus favours $1-100\xi$. No
variant acts on the lepton ratios; for $v$ and $G$ they differ only by $10^{-4}$.

\subsection{The factor~10 in \texorpdfstring{$v/E_P$}{v/E\_P}}

$v/E_P=\xi^4/(5\pi)$ with $5\pi=(\pi/2)\cdot10$ agrees to $-0.23\,\%$ at a measurement error of
$10^{-5}$. The factor~10 is fitted to the measured value: in the working state of February 2026,
$m_P/(f^4\cdot\pi/2)$ gave $2467$\,GeV instead of $246$\,GeV, and the gap of $10.02$ was closed with
$1/10$. Two readings were examined:
\begin{itemize}[leftmargin=*,itemsep=1pt]
\item Rationalised gravitation ($8\pi G=1$) with the reduced Planck mass would give
  $2\sqrt{8\pi}=10.03$ instead of~10. Docs.~374 and 376, however, show that the appropriate
  reference energy is the unreduced Planck energy; the reduced reading would be a different
  physical assumption. This reading does not hold \XM.
\item Measured versus bare $v$: if the bare $v=248.93$\,GeV of the ladder is inserted instead of
  the measured $v=246.22$\,GeV, the formula requires the denominator $15.501$. This is
  $5\pi\,K_{\text{frak}}=15.4985$ or $\pi^3/2=15.5031$, each to within $1.5\times10^{-4}$. The
  factor would then be $10\,K_{\text{frak}}\approx\pi^2$ instead of~10. The two cannot be
  distinguished at this level \SM.
\end{itemize}
The second reading fits the separation of basic form and SI: the factor~10 would then be a mix-up
of the bare value of the ladder with the SI comparison value of $v$. It is not a derivation.

% ============================================================
\section{Connections}
% ============================================================

The residuals do not stand unconnected side by side:
\begin{enumerate}[leftmargin=*,itemsep=2pt]
\item \textbf{$\alpha$ and Galois.} The $\alpha$ residual of $1.7\times10^{-4}$ is the distance of
  $m_em_\mu$ from $54$\,MeV$^2$. With the Galois value, $K_{\text{frak}}=74/75$ stays exact and the
  residual lies on the side of the masses \KM.
\item \textbf{Electron and $K_{\text{frak}}$.} The electron residual of the chain is
  $K_{\text{frak}}-1$ to within $1.6\times10^{-4}$. The chain, however, contains the factor~10;
  as long as that is open, this agreement is not robust \SM.
\item \textbf{Factor~10 and $K_{\text{frak}}$.} With the bare $v$, the factor~10 becomes
  $10\,K_{\text{frak}}$; the same correction that carries $E_0$ and $\alpha$ would then also stand
  in $v/E_P$ \SM.
\item \textbf{$\beta_T$.} The Higgs correction, at $-2.3\,\%$, is the largest and oldest. With
  FFGFT's own values for $v$ it drops to $-1.8\,\%$; it shows no robust connection with
  $K_{\text{frak}}$ \SM.
\end{enumerate}

% ============================================================
\section{Measurement precision}
% ============================================================

\begin{center}
\small
\begin{tabular}{@{}lll@{}}
\toprule
\textbf{Input} & \textbf{relative uncertainty} & \textbf{effect} \\
\midrule
$\alpha$, $m_e$ & $1.5\times10^{-10}$, $3\times10^{-10}$ & negligible \\
$m_\mu$, $M_Z$ & $2\times10^{-8}$, $2\times10^{-5}$ & negligible \\
$G$ (hence $E_P$) & $2\times10^{-5}$ ($1\times10^{-5}$) & negligible \\
$m_\tau$ & $5\times10^{-5}$ & negligible except for Koide \\
$m_h$ & $9\times10^{-4}$ & $\beta_T$ to $0.18\,\%$ \\
$m_t$, $M_W$ & $1.7\times10^{-3}$, $1.7\times10^{-4}$ & limiting \\
neutrino data, $\theta_{13}$, $\lambda_{\text{CKM}}$ & $1$ to $3\,\%$ & limiting \\
\bottomrule
\end{tabular}
\end{center}
For the lepton masses, $\alpha$, $v/E_P$ and the SI chain the measurement errors contribute only a
minimal amount; the residuals there are residuals of the theory. For $m_h$, $m_t$, $M_W$, the
neutrinos and the mixing angles the measurement is the limiting factor, and the deviations are
given there in standard deviations.

% ============================================================
\section{Conclusion}
% ============================================================

The ratios of the basic form agree to within per mille to a few per cent, and where the measurement
is precise, the residuals are genuine residuals of the theory. Three correction quantities carry
them. $\beta_T$ is the oldest: the earliest documents already set $\beta_T=1$ alongside $\alpha=1$
and saw a small deviation; today it is $0.977$, and $0.982$ with FFGFT's own $v$.
$K_{\text{frak}}=74/75$ is the correction of the SI translation; among the variants it is the only
one with a whole number of turns, and it matches the Galois form to $7.6$\,ppm. The factor~10 in
$v/E_P$ is fitted; with the bare $v$ of the ladder it becomes $10\,K_{\text{frak}}\approx\pi^2$.
Open are the cause of the lepton residuals, the residual of $\beta_T$ and the origin of the factor
in $v/E_P$.

\vspace{1em}
\noindent\textit{Check script:} \texttt{2/python/Dok387\_Skripte/pruef\_387\_abweichungen.py}
--- 27/27.
